\section{Conclusion}
\label{sec:concl}
We tested whether System~1's confidence can decide when a chess agent should search. Taking the four research questions of stage V1 in turn: (RQ1) Search lowers ACPL by \dAcplAbs\cp\ (95\% CI \dAbsLow--\dAbsHigh\cp) at \costRatio$\times$ the cost, through a few large corrections. (RQ2) Confidence is weakly informative about move quality (AUROC \aucConf) but not where search helps; blunders are more frequent at high confidence and puzzles trigger search less often than game positions. (RQ3) The confidence gate is not distinguishable from a random gate at any of four rates, and the recovery of the search benefit is proportional to the search budget; \Htwo{} and \Hthree{} are not supported. (RQ4) The hybrid is not shown to play better than System~1 alone, and its tournament estimate is slightly lower.

Two follow-up stages on fresh test months changed the agent but not the verdict on gating. A shallow material blunder filter and related repairs lowered ACPL from 114.5 to 68.7 on 2{,}000 positions and raised Elo by about 200 over System~1 alone (V2); the filter accounts for most of the ACPL reduction, but its separate effect on playing strength is not isolated, and the retrained value network made no measurable difference. A logistic gate trained to predict search gain (V3) was ahead of a random gate at every tested rate. At the primary 25\% rate the advantage is $+2.5$\cp: not significant under the pre-registered test ($[-0.9,+5.8]$, $p=0.069$; nothing survives Holm correction) and significant under a post hoc, lower-variance test ($[+0.4,+4.7]$, Holm $p=0.02$). System~1 contributes about 212 Elo over a classical policy once the filter is present. All V2 and V3 statistics were recomputed from the per-position and per-game files.

Three questions remain open: whether a gain-aware selector can capture more of the benefit that the post hoc oracle indicates (the learned gate recovers about 12\% of the gap at 25\%); whether the learned gate's modest advantage holds under a pre-specified lower-variance test on new data, with thresholds calibrated on game positions; and which of V2's bundled components, beyond the filter, matter for playing strength. In its present form the confidence-gating hypothesis is not supported. A learned gate does better than random gating, but only slightly, and a material blunder filter helped far more.
